% Generated draft: requires mathematical and source review.
\documentclass[11pt]{article}
\usepackage[margin=25mm]{geometry}
\usepackage{fontspec}
\setmainfont{Latin Modern Roman}
\usepackage{amsmath,amssymb,mathtools,unicode-math}
\setmathfont{Latin Modern Math}
\usepackage{enumitem,needspace,xurl,hyperref}
\setlength{\parindent}{0pt}
\setlength{\parskip}{5pt}
\setlength{\emergencystretch}{3em}
\newcommand{\N}{\mathbb{N}}
\newcommand{\Z}{\mathbb{Z}}
\newcommand{\Q}{\mathbb{Q}}
\newcommand{\R}{\mathbb{R}}
\newcommand{\C}{\mathbb{C}}
\newcommand{\F}{\mathbb{F}}
\newcommand{\A}{\mathbb{A}}
\newcommand{\PP}{\mathbb{P}}
\newcommand{\E}{\mathbb{E}}
\newcommand{\eps}{\varepsilon}
\newcommand{\dd}{\,\mathrm{d}}
\begin{document}
\section*{UnsolvedMath ID: 20001968}
\textbf{Source problem code:} catalogue entry
\subsection{Definitions and mathematical statement}
For a complex algebraic variety $Y$, use ordinary-cohomology mixed Hodge numbers
\[
 h^{p,q;k}(Y)=\dim\operatorname{Gr}_F^p
                       \operatorname{Gr}_{p+q}^WH^k(Y;\C),
 \qquad H(Y;x,y,t)=\sum_{p,q,k}h^{p,q;k}(Y)x^py^qt^k.
\]
This is not a compactly supported or alternating-sign convention. Write $H_n(x,y,t)$ for this polynomial of the intended twisted Betti moduli space $M_B^d(\operatorname{GL}(n,\C))$. Its precise genus, twisting and coprimality hypotheses are not included in the extracted conjecture and require the preceding source definitions.

\textbf{Conjecture 6.18, Hausel--Rodriguez--Villegas (2004).} The source proposes
\[
 H_n(\sqrt q,\sqrt q,t)=\mathbb H_n(q,t),
\]
where $\mathbb H_n$ denotes the separately constructed plethystic expression, distinguished here from the trivariate geometric polynomial. The source extraction renders both symbols as $H_n$; the defining plethystic series is absent, so the right side is not treated as an arbitrary renaming of the left side. Specifying that series and its normalization is necessary for a complete all-rank statement.

\textbf{Rank-two example as supplied.} Example 6.19, attributed to Hausel--Thaddeus (2000), gives the following reconstructed reading of the line-broken fractions:
\[
\begin{aligned}
 \frac{H_2(\sqrt q,\sqrt q,t)}{(qt+1)^{2g}}
 ={}&\frac{(q^2t^3+1)^{2g}}{(q^2t^2-1)(q^2t^4-1)}
 +\frac{q^{2g-2}t^{4g-4}(q^2t+1)^{2g}}{(q^2-1)(q^2t^2-1)}\\
 &-\frac{q^{2g-2}t^{4g-4}(qt+1)^{2g}}{2(qt^2-1)(q-1)}
 -\frac{q^{2g-2}t^{4g-4}(qt-1)^{2g}}{2(q+1)(qt^2+1)}.
\end{aligned}
\]
For $g=3$ its displayed polynomial is
\[
\begin{aligned}
 \frac{H_2(\sqrt q,\sqrt q,t)}{(qt+1)^6}
 ={}&t^{12}q^{12}+t^{12}q^{10}+6t^{11}q^{10}
      +t^{12}q^8+t^{10}q^{10}\\
 &+6t^{11}q^8+16t^{10}q^8+6t^9q^8+t^{10}q^6+t^8q^8
      +26t^9q^6\\
 &+16t^8q^6+6t^7q^6+t^8q^4+t^6q^6+6t^7q^4+16t^6q^4\\
 &+6t^5q^4+t^4q^4+t^4q^2+6t^3q^2+t^2q^2+1.
\end{aligned}
\]
The doubled plus sign in the extraction is interpreted as one plus sign. These are source examples with reconstructed-unverified status, not independently verified identities or a substitute for the missing all-rank definition.

\textbf{Supplied rank-one result and scope audit.} In rank one the closed Betti moduli space is $(\C^*)^{2g}$, and the aid asserts
\[
 H_1(x,y,t)=(1+xyt)^{2g},\qquad
 \mathbb H_1(q,t)=(1+qt)^{2g},\qquad
 h^{k,k;k}=\binom{2g}{k}
\]
for $0\le k\le2g$, with all other mixed Hodge numbers zero. Its plethystic calculation is said to use the exact AIM normalization and a corrected one-box hook length; that construction is not supplied here.

The aid also distinguishes three literature claims: HRV (2008), Corollary 4.1.11, supplies Hodge--Tateness for the closed twisted family; Mellit (2020) supplies the Poincar\'e specialization, not the full weight-graded formula; Mellit (2025), Theorem 1.5, supplies a broader generic-semisimple Hodge--Tate and curious-Lefschetz result, not the missing closed formula. These are the aid's dated assessments, not an independent current survey. Under Hodge--Tateness, $h^{p,q;k}=0$ unless $p=q$, so the diagonal $H(\sqrt q,\sqrt q,t)$ retains every coefficient $h^{p,p;k}$. Without that independent condition diagonal specialization is not injective on general trivariate polynomials; knowing a diagonal alone is not a proof of the complete mixed Hodge structure. The Poincar\'e polynomial is the further specialization $H(1,1,t)$.
\subsection{Short English statement}
The 2004 conjecture identifies a diagonal mixed-Hodge polynomial with a specific plethystic expression. The supplied aid verifies the rank-one normalization and separates general Hodge--Tateness and the Poincar\'e specialization from the still-missing full formula. The actual plethystic definition and moduli-space hypotheses are absent from this extract.
\subsection{Sources}
\begin{itemize}
\item Local UnsolvedMath registry, permanent ID 20001968, source code catalogue entry; canonical \texttt{aim-geometry-notes.json}, notes[305], statement status reconstructed unverified.
\item AIM 2004 workshop \emph{Moment maps and surjectivity in various geometries}, Conjecture 6.18 (Hausel--Rodriguez--Villegas 2004), Example 6.19 (Hausel--Thaddeus 2000): \url{https://aimath.org/WWN/momentmaps/momentmaps.pdf}.
\item Research-summary references as supplied: HRV 2008, Corollary 4.1.11; Mellit 2020; Mellit 2025, Theorem 1.5. No full titles, URLs or bibliographic entries for these works are supplied.
\item Registry research attempt 1, dated 2026-08-14: partial result, rigor ``proved'', novelty confidence low. The explicitly machine-generated, non-peer-reviewed aid requires independent verification before citation.
\end{itemize}
\end{document}
